\paragraph{A∞-模。}\label{sub:a-infinity-modules}
设 $Ch$ 是由 $\KK$-\trVectorSpace{}的\trChainComplex{}组成的 dg \trCategory，视作一个%
\trACategory。取定一个\trNonunital{} \trACategory{} $\scr{A}$，把 $\scr{A}$ 上%
\textbf{\trNonunital{}（右）\trAinfty-模}（non-unital (right) $A_\infty$-modules over
$\scr{A}$）的%
\trNonunital{} \trACategory{}定义为
\begin{equation*}
\mathit{nu}\text{-}\mathit{mod}(\scr{A})=\mathit{nu}\text{-}\mathit{fun}(\scr{A}^{\mathrm{opp}},Ch).
\end{equation*}
值得把细节具体地铺陈出来。$\scr{Q}=\mathit{nu}\text{-}\mathit{mod}(\scr{A})$ 的一个%
\trObject{} $\scr{M}$ 由如下数据组成：对每个 $X\in\Ob\scr{A}$，有一个\trGraded\trVectorSpace{}%
$\scr{M}(X)$，连同各阶 $d\ge 1$ 的映射
\begin{equation*}
\mu^d_{\scr{M}}\colon\scr{M}(X_{d-1})\otimes\hom_{\scr{A}}(X_{d-2},X_{d-1})
\otimes\cdots\otimes\hom_{\scr{A}}(X_0,X_1)
\longrightarrow\scr{M}(X_0)[2-d].
\end{equation*}
若按 \trAFunctor{}的语言来想，则 $\mu^1_{\scr{M}}$ 就是像\trObject{}%
$\scr{M}(X_0)\in\Ob Ch$ 上的\trDifferential，而当 $d\ge 2$ 时，$\mu^d_{\scr{M}}$ 是%
\trFunctor{}在\trMorphism{}上的作用，更确切地说即 $\scr{M}^{d-1}$（把诸自变量的次序
倒换）。此时\eqref{eq:afunctor-equations} 就成为
\begin{equation}\label{eq:module-associativity-equation}
\begin{split}
&\sum_n(-1)^{\clubsuit_n}\,\mu^{n+1}_{\scr{M}}
\bigl(\mu^{d-n}_{\scr{M}}(b,a_{d-1},\ldots,a_{n+1}),\ldots,a_1\bigr)\\
&\quad+\sum_{m,n}(-1)^{\clubsuit_n}\,\mu^{d-m+1}_{\scr{M}}
\bigl(b,a_{d-1},\ldots,a_{n+m+1},\mu^m_{\scr{A}}(a_{n+m},\ldots,a_{n+1}),\ldots,a_1\bigr)=0,
\end{split}
\end{equation}
其中第二个求和取遍 $n+m<d$。接下来，设 $\scr{M}_0$、$\scr{M}_1$ 为两个%
\trNonunital{} \trAinfty-模。
$\hom_{\scr{Q}}(\scr{M}_0,\scr{M}_1)$ 中的\trMorphism{}称为%
\textbf{预\trModule\trHomomorphism}（pre-module homomorphism），它们写成一系列映射
\begin{equation*}
t^d\colon\scr{M}_0(X_{d-1})\otimes\hom_{\scr{A}}(X_{d-2},X_{d-1})
\otimes\cdots\otimes\hom_{\scr{A}}(X_0,X_1)
\longrightarrow\scr{M}_1(X_0)[|t|-d+1],
\end{equation*}
$d\ge 1$。这与预\trNaturalTransformation{} $T$ 的一般语言之间的联系，只在于自变量的
明显对调，外加一个投机取巧的正负号改动：
\begin{equation}\label{eq:premodule-sign-convention}
t^d(b,a_{d-1},\ldots,a_1)=(-1)^{\text{\S}}\,T^{d-1}(a_1,\ldots,a_{d-1})(b),
\end{equation}
其中 $T^{d-1}(a_1,\ldots,a_{d-1})\in\hom_{Ch}(\scr{M}_0(X_{d-1}),\scr{M}_1(X_0))$，
而 $\text{\S}=(|T|-1)|b|+|T|(|T|-1)/2$。真正的\trModule\trHomomorphism{}是那些使
$\mu^1_{\scr{Q}}(t)=0$ 的 $t$。为了看出这一点的具体含义，我们把 $\scr{Q}$ 中的%
\trComposition{}映射写下来，它们是：
\begin{equation}\label{eq:module-composition-maps}
\begin{split}
(\mu^1_{\scr{Q}}t)^d(b,a_{d-1},\ldots,a_1)
&=\sum_n(-1)^{\ddagger}\,\mu^{n+1}_{\scr{M}_1}
\bigl(t^{d-n}(b,a_{d-1},\ldots,a_{n+1}),a_n,\ldots,a_1\bigr)\\
&\quad+\sum_n(-1)^{\ddagger}\,t^{n+1}
\bigl(\mu^{d-n}_{\scr{M}_0}(b,a_{d-1},\ldots,a_{n+1}),a_n,\ldots,a_1\bigr)\\
&\quad+\sum_{m,n}(-1)^{\ddagger}\,t^{d-m+1}\bigl(b,a_{d-1},\ldots,\\
&\qquad\mu^n_{\scr{A}}(a_{n+m},\ldots,a_{n+1}),a_n,\ldots,a_1\bigr);
\end{split}
\end{equation}
\begin{equation*}
(\mu^2_{\scr{Q}}(t_2,t_1))^d(b,a_{d-1},\ldots,a_1)
=\sum_n(-1)^{\ddagger}\,t_2^{n+1}
\bigl(t_1^{d-n}(b,a_{d-1},\ldots,a_{n+1}),a_n,\ldots,a_1\bigr);
\end{equation*}
\begin{equation*}
\mu^d_{\scr{Q}}=0\quad\text{对一切 }d\ge 3\text{，}
\end{equation*}
其中 $\ddagger=|a_{n+1}|+\cdots+|a_{d-1}|+|b|-d+n+1$。$\mu^3_{\scr{Q}}$、$\mu^4_{\scr{Q}}$
等等的消没，是下述非正式却有用的普遍原理的一个实例：$\mathit{nu}\text{-}\mathit{fun}(\scr{A},\scr{B})$
继承的是 $\scr{B}$ 的性质，而不是 $\scr{A}$ 的性质。

在\trCohomology{}层面上，任何 $\scr{M}$ 都给出一个\trNonunital{}%
$H(\scr{A})$-\trModule{} $H(\scr{M})$；说得更形式化一点，它就是一个从 $H(\scr{A})$ 到%
\trGraded\trVectorSpace{}的\trCategory{}的\trNonunital\trContravariant\trFunctor。它由 $\scr{M}(X)$
关于 $\partial(b)=(-1)^{|b|}\mu^1_{\scr{M}}(b)$ 的\trCohomology，以及由
$b\cdot a=(-1)^{|a|}\mu^2_{\scr{M}}(b,a)$ 诱导的\trModule{}结构组成。类似地，给定%
\trAinfty-\trModule\trHomomorphism{} $t$，我们得到诱导映射
$H(t)\colon H(\scr{M}_0(X))\to H(\scr{M}_1(X))$，$[b]\mapsto[(-1)^{|b|}t^1(b)]$，它们在
通常意义下构成一个\trModule\trHomomorphism。下面这条是\cref{thm:right-composition-iso}
的特例：
\vkthmsync{16}
\begin{lemma}\label{thm:module-composition-quasiisomorphism}
设 $t\colon\scr{M}_0\to\scr{M}_1$ 是一个\trAinfty-\trModule\trHomomorphism，使得对每个
$X\in\Ob\scr{A}$，映射 $H(t)\colon H(\scr{M}_0(X))\to H(\scr{M}_1(X))$ 都是%
\trIsomorphism。那么，对任意 $\scr{N}$，用 $t$ 作左\trComposition{}与右\trComposition{}%
分别诱导\trQuasiisomorphism{}%
$\hom_{\scr{Q}}(\scr{M}_1,\scr{N})\to\hom_{\scr{Q}}(\scr{M}_0,\scr{N})$ 与
$\hom_{\scr{Q}}(\scr{N},\scr{M}_0)\to\hom_{\scr{Q}}(\scr{N},\scr{M}_1)$。
\noqed
\end{lemma}
特别地，若 $\scr{M}$ 是满足 $H(\scr{M}(X),\mu^1_{\scr{M}})=0$（对一切 $X$）的\trNonunital{}
\trAinfty-模，则对一切 $\scr{N}$ 都有
$H(\hom_{\scr{Q}}(\scr{M},\scr{N}))=H(\hom_{\scr{Q}}(\scr{N},\scr{M}))=0$。
借用 [136]、[21] 的术语来说：在 \trAinfty-模的世界里，万事万物既是\textbf{$K$-投射}
（$K$-projective）的，又是\textbf{$K$-内射}（$K$-injective）的。
